\section{Disaster Recovery}

\subsection{Critical Dependencies Chain}

The lab infrastructure has a cascading dependency chain. A failure in one layer propagates:

\begin{enumerate}
    \item \textbf{Time (NTP)} $\rightarrow$ if time drift exceeds 128 seconds, Kerberos tickets expire and TLS handshakes fail
    
    \item \textbf{DNS (ns1/BIND)} $\rightarrow$ without DNS resolution, all hostname-based services (Kerberos SRV records, NTP upstream, LDAP) break
    
    \item \textbf{Kerberos (odroid-c1/KDC)} $\rightarrow$ without the KDC, authentication fails for all Kerberos-authenticated services
    
    \item \textbf{LDAP (bbb1/slapd)} $\rightarrow$ without LDAP, directory lookups fail and user/service account information is unavailable
\end{enumerate}

Recovery priority should follow this chain in reverse: fix time first, then DNS, then KDC, then LDAP.

\subsection{Failure Scenarios and Recovery}

\subsubsection{NTP Failure -- All Hosts Lose Time Sync}

Symptoms: \texttt{timedatectl} shows "System clock synchronized: no", Kerberos tickets rejected.

Recovery:
\begin{enumerate}
    \item Verify GPS unit is locked on odroid-c1 (time host)
    
    \item Check ntpd/ntpsec service: \texttt{systemctl status ntpsec}
    
    \item If GPS is offline, enable holdover mode (internal oscillator continues at reduced accuracy)
    
    \item On affected clients, force resync: \texttt{ntpdate -s time.lab.bitsmasher.net} or wait for ntpd's step-sync to kick in
    
    \item Verify with \texttt{ntpq -p} on clients and \texttt{ntptime} on the server
\end{enumerate}

\subsubsection{DNS (ns1) Offline -- Cascade Failure}

Symptoms: hostnames don't resolve, Kerberos SRV discovery fails.

Recovery:
\begin{enumerate}
    \item Check ns1 reachability: ping from any known-good host
    
    \item If ns1 is up but not answering: check BIND (named) service status
    
    \item If BIND has crashed due to config error: restore from last known-good named.conf
    
    \item As a workaround, add manual entries to /etc/hosts on affected hosts until DNS recovers
    
    \item Update krb5.conf with explicit \texttt{kdc = odroid-c1.lab.bitsmasher.net}
\end{enumerate}

\subsubsection{LDAP (bbb1) Down -- No Directory Service}

Symptoms: slapd failed, status.sh reports both slapd and ldapsearch as failing.

Recovery:
\begin{enumerate}
    \item SSH to bbb1/lab.bitsmasher.net as root
    
    \item Check TLS certificate validity in the configured cert path
    
    \item Regenerate or replace expired certificates
    
    \item Restart slapd: \texttt{systemctl restart slapd}
    
    \item Verify with ldapsearch for franklin and sly DNs
\end{enumerate}

\subsubsection{KDC (odroid-c1) Offline -- No Authentication}

Symptoms: kinit fails for all realms, service authentication denied.

Recovery:
\begin{enumerate}
    \item SSH to odroid-c1 as root (franklin user's key is not authorized)
    
    \item Check kdc process: \texttt{kadmind/krb5kdc status}
    
    \item Verify KDC database integrity: \texttt{kdb5_util list}
    
    \item Restart kerberos services if needed
    
    \item Regenerate host keytabs for affected services via \texttt{kadmin.local}
\end{enumerate}

\subsection{Backup and Restoration}

\begin{itemize}
    \item Ansible playbooks (lab-franklin collection) are the single source of truth -- they can reconstruct any configured state
    
    \item Bare git repos with GPG-encrypted pass entries handle credential backup
    
    \item KDC database (\texttt{\$KRB5\_KDB\_FILE}) should be backed up regularly on odroid-c1
    
    \item DNS zone files should have local copies (in the DNS role's files/ directory)
\end{itemize}

\subsection{Prevention: Monitoring}

The current monitoring approach is manual via status.sh. For improved resilience, consider:
\begin{itemize}
    \item Automated ping/SSH checks on critical hosts in HEARTBEAT.md
    
    \item Nagios/Zabbix-style monitoring on a dedicated host
    
    \item Email/slack alerts for service down conditions
\end{itemize}

The lab-franklin Ansible collection should eventually include a monitoring role to automate this.
